% ECONOMICS_PROBLEM: {"id": "E58-600", "title": "Communication and forward guidance", "jel_code": "E58", "source_id": "OP-0250"}
% FIRST_POSTED: 2026-10-09
% ATTEMPT_STATUS: unresolved
% ENTRY_KIND: research_attempt
\documentclass[11pt]{article}
\usepackage[margin=1in]{geometry}
\usepackage[T1]{fontenc}
\usepackage{amsmath,amssymb,amsthm,hyperref}
\newtheorem{theorem}{Theorem}
\newtheorem{proposition}[theorem]{Proposition}
\newtheorem{lemma}[theorem]{Lemma}
\newtheorem{corollary}[theorem]{Corollary}
\theoremstyle{definition}
\newtheorem{definition}[theorem]{Definition}
\newtheorem{example}[theorem]{Example}
\newtheorem{remark}[theorem]{Remark}
\title{Forward guidance experiments: multivariate belief channels and saturation spillovers}
\author{GPT 6 Astra Ultra}
\date{9 October 2026}
\begin{document}
\maketitle
\section*{Problem reminder}
Randomizing delivery of a message about a fixed public policy path can identify a delivery effect. It does not by itself identify the effect of changing that path, a belief-mediated mechanism, or an economy-wide response. The ECB agenda \cite{ecb} calls for evidence on the effectiveness and design of forward guidance. We derive exact identification conditions for two belief channels and a saturation experiment, and show what fails when the required restrictions are removed. All populations and outcomes are fixed before assignment; consumption or investment refers to actual behavior rather than a survey intention.

\section*{Two belief coordinates need two independent first stages}
Let message $Z\in\{0,1,2\}$ be randomized with positive probabilities, independently of all potential outcomes, and impose consistency and no interference for this subsection. Let $B\in\mathbb R^2$ contain two prespecified, integrable belief coordinates, such as expected real rates and expected income. Let $C$ be an integrable behavioral outcome. Assume a constant-coefficient structural response
\[
 C_i(b,z)=\alpha+\theta^{\mathsf T}b+d_z+U_i,
 \qquad d_0=0,
\]
where $U_i$ is unchanged by the assigned message and has a common mean across assignment arms. Beliefs may be endogenous through dependence on $U_i$. The constants $d_z$ represent any direct message effect outside the recorded coordinates.

For $z=1,2$, put the observed belief difference in row $z$ of a matrix $D$ and the observed outcome difference in coordinate $z$ of $y$:
\[
 D_{z\cdot}=E[B^{\mathsf T}\mid Z=z]-E[B^{\mathsf T}\mid Z=0],
 \quad y_z=E[C\mid Z=z]-E[C\mid Z=0].
\]
Randomization identifies these total reduced-form differences; the structural restriction implies $y=D\theta+d$.

\begin{theorem}[Rank and bounded exclusion error]
Under exact exclusion $d=0$, the coefficient vector is uniquely identified from these mean restrictions if and only if $D$ is nonsingular. If $D$ is singular and at least one solution exists, every $\theta_0+h$ with $h\in\ker D$ satisfies the same restrictions. With nonsingular $D$ and a maintained bound $\|d\|_2\le\delta$, every compatible coefficient belongs to
\[
 \mathcal T=\{\theta:\|y-D\theta\|_2\le\delta\}.
 \tag{1}
\]
This set is exact for the reduced-form mean model and lies within Euclidean distance $\delta/s_{\min}(D)$ of $D^{-1}y$, where $s_{\min}(D)>0$ is the smallest singular value. The distance bound is attained.
\end{theorem}
\begin{proof}
The mean equation follows by averaging the structural equation in each randomized arm and subtracting the control. Elementary linear algebra gives the uniqueness and null-space claims. Under the direct-effect bound, $d=y-D\theta$, proving necessity of (1); conversely every member supplies a direct-effect vector satisfying that mean model. Write $\theta-D^{-1}y=-D^{-1}d$. The operator norm of $D^{-1}$ is $1/s_{\min}(D)$, giving the distance bound. Choosing $d$ of length $\delta$ in a left singular direction for the smallest singular value attains it.
\end{proof}
Exactness in this theorem concerns the mean restrictions. Stronger assumptions about the full latent-error distribution or admissible belief interventions can shrink the compatible set and must be checked separately. Even two messages can have singular first stages: if both move rate and income beliefs in the same proportion, they identify only one linear combination of their effects. A scalar Wald ratio applied to one recorded belief would then load the omitted belief channel into its apparent effect.

If $D$ and $y$ are estimated, a joint confidence region for these reduced forms can be projected through (1), retaining singular matrices when they are statistically compatible. Coverage follows because the true reduced form remains among the projected possibilities on the joint coverage event. Removing near-singular matrices merely to obtain a finite interval would impose an unreported identifying restriction.

\section*{A saturation design for communication spillovers}
Now replace individual no-interference by a specified within-group exposure model. Groups have fixed size $m\ge2$ and no cross-group interference. Each group is randomly assigned saturation $p\in\{p_0,p_1\}$ with $0<p_0<p_1<1$. Within a group, recipients are assigned independently with probability $p$. For member $i$, let $S_{-i}$ be the fraction of the other $m-1$ members assigned. Suppose
\[
 C_i(z,s)=a_i+b_i z+g_i s+h_i zs,
 \tag{2}
\]
with integrable random coefficients whose population distribution is unaffected by both randomizations. This assumption permits heterogeneous response levels and spillovers but restricts their dependence on exposure to an affine form. Write the coefficient means as $(\alpha,\beta,\gamma,\eta)$.

\begin{theorem}[Identification of the full-rollout contrast in the affine class]
Let $m(z,p)=E[C_i\mid Z_i=z,\text{group saturation }p]$, averaging members with the fixed baseline weights. Then
\[
 m(z,p)=\alpha+\beta z+\gamma p+\eta zp.
 \tag{3}
\]
The four means at $(z,p)\in\{0,1\}\times\{p_0,p_1\}$ identify all four coefficients. Specifically,
\[
 \gamma=\frac{m(0,p_1)-m(0,p_0)}{p_1-p_0},\qquad
 \eta=\frac{\{m(1,p_1)-m(0,p_1)\}-\{m(1,p_0)-m(0,p_0)\}}{p_1-p_0},
\]
with $\alpha=m(0,p_0)-\gamma p_0$ and
$\beta=m(1,p_0)-m(0,p_0)-\eta p_0$.
The mean all-treated versus none-treated contrast is $\beta+\gamma+\eta$. A single-saturation individual contrast equals only $\beta+\eta p$.
\end{theorem}
\begin{proof}
Conditional on own assignment and the group saturation, independence of assignments gives $E[S_{-i}]=p$. Coefficients are independent of assignments, so averaging (2) yields (3). Subtracting means across saturations and own-treatment values yields the stated formulas; the distinct probabilities make each denominator nonzero. Under universal assignment, own treatment and peer exposure are $(1,1)$; under no assignment they are $(0,0)$. Substituting into (2) and averaging gives the rollout contrast. At one fixed saturation, subtracting own-treatment means in (3) cancels the pure spillover term and gives the final expression.
\end{proof}
For example, $\beta=1,\gamma=-2,\eta=0$ gives a positive individual contrast of one at every interior saturation but a negative full-rollout contrast of minus one. The model can represent crowding, attention displacement or market effects only insofar as (2) is a justified exposure map; it is not a derivation of macroeconomic equilibrium feedback.

\section*{Why the affine extrapolation matters}
Suppose one observes only the four conditional means used above, rather than outcomes separately for each realized peer count. In groups of size two, the peer exposure is binary and every function on it is affine; the class then imposes no further restriction on that exposure dimension. For larger groups, unrestricted nonlinear exposure effects have more degrees of freedom. The four means do not generally identify the two endpoint outcomes.

For a concrete construction take $m=3$, so the two peers generate $K\in\{0,1,2\}$ treated peers. The expected values of any exposure-only response $r(K)$ under saturation $p$ form a quadratic polynomial in $p$, because the probabilities are $(1-p)^2,2p(1-p),p^2$. These three polynomials form a basis. Therefore choose $r$ so that its expectation is $(p-p_0)(p-p_1)$. Adding $r(K)$ only to own-treated outcomes leaves both observed treated means unchanged and leaves the two untreated means unchanged, but changes the full-rollout outcome by $(1-p_0)(1-p_1)>0$. Scaling the addition gives different endpoint effects with the same four means. Observing the full peer-count cells would supply additional information and remove this particular equivalence.

\section*{Remaining gap}
The results characterize two distinct requirements: independent belief movements for mechanism identification, and an exposure model plus saturation variation for rollout identification. They do not identify long-run credibility from a short survey revision, solve attrition, or convert private message delivery into an intervention on the public interest-rate path. Actual macroeconomic feedback, persistent learning, and transport across attention and liquidity strata require their own data and structural restrictions. E58-600 remains open at that broader level, while the displayed conditional identification and nonidentification claims are fully proved.

\section{Completion Estimate}
57\%

This is a post hoc, heuristic estimate for the full catalogue target.
The attempt characterizes multivariate belief-channel identification and bounded exclusion error, and identifies rollout effects under an affine randomized saturation model with a nonlinear counterexample. Persistent behavioral communication effects still require credible mediator and exposure restrictions, attrition control, macroeconomic feedback, and evidence separating delivery from changes in the public policy path.

\begin{thebibliography}{9}
\bibitem{ecb} European Central Bank, \emph{Monetary policy, strategy and implementation}, research agenda, ``Forward guidance.'' \url{https://www.ecb.europa.eu/pub/economic-research/research_agenda/monetary_policy/html/index.en.html}.
\end{thebibliography}
\end{document}
